\documentclass[11pt]{article}
\usepackage[T1]{fontenc}
\usepackage{lmodern,amsmath,amssymb,amsthm,microtype}
\usepackage[margin=26mm,top=23mm,bottom=24mm]{geometry}
\usepackage[hidelinks]{hyperref}
\usepackage{fancyhdr}
\pagestyle{fancy}\fancyhf{}\lhead{\small Lacunary dilates}\rhead{\small JSP-000828}\cfoot{\thepage}
\setlength{\headheight}{14pt}
\newtheorem*{theorem}{Theorem}
\newcommand{\T}{\mathbb T}
\newcommand{\E}{\mathbb E}
\newcommand{\1}{\mathbf 1}
\newcommand{\card}{\#}
\begin{document}
\begin{center}
{\huge\bfseries The sharp almost-everywhere growth\\[3pt] of sums of lacunary dilates\par}
\vspace{9pt}
{\large A self-contained proof for Erd\H{o}s problem 995}\\[4pt]
21 November 2026
\end{center}

\section*{1. The answer}
On the circle $\T=\mathbb R/\mathbb Z$ with Lebesgue probability measure $\mu$, write
\[
 S_N(f,x)=\sum_{k=1}^N f(n_kx).
\]
The integers $n_k$ are positive. A sequence is lacunary if
$n_{k+1}/n_k\ge q$ for some $q>1$ and every $k$. For every such sequence and
every $f\in L^2(\T)$,
\[
                    S_N(f,x)=o\!\left(N\sqrt{\log N}\right)
                    \quad\text{for almost every }x.                 \tag{1}
\]
For each positive sequence $M_N=o(N\sqrt{\log N})$, there exist a lacunary
sequence and a real, mean-zero $f\in L^2(\T)$ such that
\[
            \limsup_{N\to\infty}\frac{S_N(f,x)}{M_N}=+\infty
                    \quad\text{for almost every }x.                 \tag{2}
\]
Thus the proposed $o(N\sqrt{\log\log N})$ bound fails, even with mean zero.
The example may depend on the proposed normalizer. The exceptional null set
in (1) may depend on $f$ and the sequence.

Here is the exact criterion, including normalizers that oscillate.
For finite $M_N>0$ put
\[
 D_j=[2^j,2^{j+1})\cap\mathbb N,\qquad
 w_j=\max_{N\in D_j}\frac{N}{M_N},\qquad
 K=\sup_{u>0}u\sqrt{\card\{j\ge0:w_j>u\}}.
\]
\begin{theorem}
If $K<\infty$, then $S_N(|f|,x)/M_N\to0$ almost everywhere for every
$f\in L^2(\T)$ and every sequence of positive integers. If $K=\infty$,
then (2) holds for a mean-zero $f\in L^2(\T)$ and a sequence of
powers of two with successive ratios at least two.
\end{theorem}

\paragraph{Relation to earlier work.}
The sharp scale follows from the general results of Quas and Wierdl
\cite{QW}. Their Theorem 2.3 applies to averages of measure-preserving
maps; their Theorem B(2) gives sharpness for iterates of the doubling map.
Their weak sequence norm leads to the criterion above. Ho's separate
construction \cite{Ho} disproves the proposed bound and gives estimates
up to an arbitrarily small logarithmic exponent. We give direct proofs of
both implications for circle dilates, including an elementary construction
of all the required counterexamples.

\newpage
\section*{2. The universal upper estimate}
For nonnegative $g$ put
\[
 B_jg(x)=2^{-j-1}\sum_{k=1}^{2^{j+1}}g(n_kx).
\]
The map $x\mapsto nx\bmod 1$ preserves $\mu$ for every positive integer $n$,
as seen by splitting $[0,1)$ into $n$ intervals of length $1/n$.
Consequently $B_j$ preserves constants and integrals. For $N\in D_j$,
positivity and $2^{j+1}/N\le 2$ give
\[
                  \frac{S_N(|f|,x)}{M_N}\le2w_jB_j|f|(x).             \tag{3}
\]
Fix $u>0$. With $h=|f|$, the part of $h$ at most $u/(2w_j)$ contributes
at most $u/2$ to $w_j B_jh$. Markov's inequality on the remaining part
therefore yields
\[
       \mu\{w_jB_jh>u\}\le\frac{2w_j}{u}
                    \int_{\{h>u/(2w_j)\}}h\,d\mu.                    \tag{4}
\]
When $K<\infty$, the number of $j$ with $w_j>s$ is at most $K^2/s^2$.
Integrating this counting bound gives, for $a>0$,
\[
 \sum_{j:w_j>a} w_j
 =a\card\{j:w_j>a\}+\int_a^\infty\card\{j:w_j>s\}\,ds
 \le\frac{2K^2}{a}.
\]
Summing (4), interchanging the nonnegative sum and integral, and
substituting $a=u/(2h(x))$ where $h(x)>0$, we obtain
\[
             \sum_j\mu\{w_jB_jh>u\}\le \frac{8K^2}{u^2}\int h^2\,d\mu .
\]
The first Borel--Cantelli lemma, applied to $u=1,1/2,1/3,\ldots$,
gives $w_jB_jh\to0$ almost everywhere. Estimate (3) proves convergence
along all $N$. Replacing an $L^2$ representative on a null set changes
its dilates only on a countable union of null sets.

\paragraph{How the criterion gives the endpoint.}
For $M_N=N\sqrt{1+\log N}$ we have $w_j=(1+j\log2)^{-1/2}$ and $K<\infty$.
Conversely, if $M_N=o(N\sqrt{\log N})$, evaluation at $N=2^j$ gives
$w_j\sqrt j\to\infty$. For every $A>0$ and all sufficiently large $J$,
the $J$ indices in $[J,2J)$ then have $w_j>A/\sqrt{2J}$. Hence
$K\ge A/\sqrt2$ for every $A$. It remains to prove (2) when $K=\infty$.

\newpage
\section*{3. Constructing the counterexample}
Ignore the null set of ambiguous binary expansions. The binary digits
$\xi_1(x),\xi_2(x),\ldots$ of a uniformly distributed $x$ are independent
fair bits: any prescribed string of $s$ initial bits has measure $2^{-s}$.
Multiplication by $2^d$ shifts these digits to the left by $d$ places.

\paragraph{Choice of sample indices.}
Let $v_j=\min(w_j,1)$. Its weak $\ell^2$ norm is still infinite: truncation
either changes finitely many terms or leaves infinitely many terms equal
to one. Removing any finite initial segment preserves this property.
At stage $m\ge1$, choose $0<\alpha<1$ and a finite set of $H$ indices
$j$ of one parity, beyond all previous indices by at least two, such that
\[
                v_j>\alpha,\qquad \alpha\sqrt H\ \ge\ 2^{m+4}m.       \tag{5}
\]
This is possible by the definition of the weak norm, retaining the more
frequent parity. For each chosen $j$ take a maximizer $N\in D_j$ of
$N/M_N$. In the increasing list of selected $N$, including earlier
stages, each new $N$ is more than twice its predecessor. The sets of
sample indices
\[
                     I_N=\{\lfloor N/2\rfloor+1,\ldots,N\}
\]
are therefore disjoint and satisfy $|I_N|\ge N/2$.

\paragraph{A rare event with many opportunities.}
Within this stage choose an integer $L\ge 2\max N$. Put $r=\lceil\log_2(LH)\rceil$,
so that $1/(2H)<L2^{-r}\le 1/H$. Let $E_m$ be the set where the
word $1\,0^{r-1}$ starts at one of the digit positions $1,\ldots,L$.
The union bound gives $\mu(E_m)\le 1/H$. Set
\[
                         g_m=(8m/\alpha)\1_{E_m}.
\]
Equation (5) implies $\|g_m\|_2\le 2^{-m}$.

Assign strictly increasing nonnegative integer exponents $a_k$ in index
order. On each $I_N$ assign consecutive exponents $d,\ldots,d+s-1$,
where $s=|I_N|$. Choose the initial $d$ large enough that its digit
interval $[d+1,d+L+r]$ is disjoint from and to the right of every
such interval previously assigned. The finitely many intervening sample
indices can be assigned consecutive exponents first, so there is no
obstruction. The stage maxima grow without bound, which assigns every $k$.

If the word starts anywhere from position $d+s$ to $d+L$, then
$2^{a_k}x\in E_m$ for every $k\in I_N$: subtracting any such $a_k$
puts the start in $[1,L]$. Call this event $C_N$. There are
$h=L-s+1\ge L/2$ possible starts. Put $\lambda=h2^{-r}$, so
$1/(4H)<\lambda\le 1/H$. Two distinct starts less than $r$ apart are
incompatible because the word starts with one and all its other bits
are zero. Starts at least $r$ apart use disjoint digits. The two-term
inclusion--exclusion bound now gives
\[
               \mu(C_N)\ge h2^{-r}-\tbinom h2 2^{-2r}
                          \ge\lambda-\lambda^2/2\ge 1/(8H).          \tag{6}
\]

\newpage
\section*{4. From the digit events to divergence}
All the $C_N$ are independent because the digit intervals used for
different selected $N$ are disjoint, including across stages. Thus
the chance of at least one $C_N$ in stage $m$ is at least
\[
              1-(1-1/(8H))^H\ge 1-e^{-1/8}.
\]
These stage events are themselves independent. The second
Borel--Cantelli lemma implies that infinitely many stages succeed
for almost every $x$.

Set $F=\sum_{m\ge1}g_m$. The summands are nonnegative and their
$L^2$ norms are summable. By the triangle inequality on partial sums
and monotone convergence of their squares, $F\in L^2$ and is finite
almost everywhere. Write $b=\int F\,d\mu$ and put $f=F-b$,
choosing any finite representative on the exceptional null set.
Discard also its countably many null preimages under the selected
dilations. On a successful stage, choose $N$ with $C_N$. For this $N$,
the sample sequence $n_k=2^{a_k}$ satisfies
\[
 \frac{S_N(F,x)}{N}\ \ge\ \frac{|I_N|}{N}\frac{8m}{\alpha}
                  \ \ge\ \frac{4m}{\alpha}.
\]
Since $\alpha<1$, for $m>b$ it follows that
\[
 \frac{S_N(f,x)}{M_N}
   =\frac{N}{M_N}\left(\frac{S_N(F,x)}{N}-b\right)
   \ge\alpha\left(\frac{4m}{\alpha}-b\right)\ge 3m .
\]
This proves (2). Strict increase of the integer exponents gives
$n_{k+1}/n_k\ge2$. Taking $M_N=N\sqrt{\log\log N}$ for all sufficiently
large $N$ supplies the counterexample in the problem.

\paragraph{A fixed sequence from the literature.}
For a shorter route to sharpness against every strictly smaller order,
use Theorem B(2) of \cite{QW} with $p=2$, $T(x)=2x\bmod1$ and
iterate times $2^k$. The map is measure-preserving and aperiodic:
for each $a\ge1$ its points of period dividing $a$ satisfy
$(2^a-1)x=0\bmod1$, a finite set. Thus the \emph{integer dilates}
in this application are $2^{2^k}$. For a centered example and a tiny
$M_N$, apply the theorem to $\max(M_N,N)$ and take the absolute value
of its witness before centering. The normalized uncentered limsup is
infinite, and the cost of centering at that larger denominator is bounded.

\begin{thebibliography}{3}\small
\bibitem{E} P. Erd\H{o}s, \emph{Problems and results on diophantine approximations},
Compositio Math. 16 (1964), 52--65, especially p.~58.
\href{https://www.numdam.org/item/CM_1964__16__52_0/}{Original article}.
\href{https://www.erdosproblems.com/995}{Modern problem statement}.
\bibitem{QW} A. Quas and M. Wierdl, \emph{Rates of divergence of non-conventional
ergodic averages}, Ergodic Theory Dynam. Systems 30 (2010), 233--262.
\href{https://doi.org/10.1017/S0143385709000054}{Publisher}.
\href{https://arxiv.org/pdf/math/0409001}{Preprint}, Thm.~B, Thms.~2.3--2.4 and 6.2.
\bibitem{Ho} B. S. Ho, \emph{Counterexamples for lacunary dilates via dyadic
spike blocks}, \href{https://arxiv.org/html/2604.18535v3}{arXiv:2604.18535v3}
(23 October 2026), finite-$L^p$ growth in Sections 6 and 8.
\end{thebibliography}
\end{document}
